\chapter{Modello Relazionale e Progettazione Logica}

\section{Introduzione}

Il passaggio dalla progettazione concettuale alla progettazione logica rappresenta la transizione dal \textit{cosa} deve essere rappresentato al \textit{come} organizzare concretamente i dati all'interno di un sistema di gestione di basi di dati (DBMS). Questo processo richiede l'adozione di uno specifico \textbf{modello dei dati}.

Sebbene la storia dell'informatica abbia visto l'evoluzione di diversi paradigmi — come il modello gerarchico, quello reticolare e quello orientato agli oggetti —, la nostra trattazione si concentrerà sul \textbf{modello relazionale}, che costituisce lo standard industriale prevalente. Vale la pena notare che i DBMS moderni più diffusi (come PostgreSQL e Oracle) adottano un'evoluzione di questo approccio, denominata \textbf{modello relazionale a oggetti}, la quale estende le fondamenta del modello relazionale puro integrando costrutti tipici del paradigma OOP.

\medskip

Uno dei principali vantaggi nell'aver adottato i \textbf{Class Diagrams UML} nella fase concettuale risiede nella loro assoluta indipendenza tecnologica: essi si adattano teoricamente a qualunque modello logico finale. Di contro, questa spiccata espressività introduce un inevitabile \textbf{disallineamento semantico} (\textit{impedance mismatch}) quando l'obiettivo è la traduzione in uno schema relazionale puro. 

Nel passaggio alla progettazione logica relazionale, le principali sfide strutturali da affrontare saranno dettate dai rigidi vincoli di atomicità e piattezza del dato tipici delle tabelle:
\begin{itemize}
    \item \textbf{Attributi strutturati:} Il modello relazionale esige che i valori degli attributi siano atomici e non strutturati. Non è quindi possibile nidificare record, oggetti o tipi di dato complessi all'interno di una colonna.
    \item \textbf{Attributi a valore multiplo (multivalore):} Ogni cella di una tabella relazionale può contenere esclusivamente un singolo valore atomico, vietando l'uso nativo di collezioni, array, liste o set all'interno di un singolo attributo.
    \item \textbf{Assenza di ereditarietà:} Il modello relazionale puro non supporta nativamente le gerarchie di generalizzazione-specializzazione, costringendo il progettista ad applicare strategie di ``appiattimento'' o di scomposizione strutturale.
\end{itemize}


\section{Concetti Fondamentali}

Il \textbf{modello relazionale} si basa sul concetto matematico di relazione. Questo modello rappresenta la struttura logica e piatta su cui verranno mappati i dati dell'applicazione. Per formalizzare rigorosamente il modello, è necessario definire i suoi tre componenti atomici: i domini, lo schema e la relazione stessa.

\medskip

Un \textbf{dominio} (indicato con $D$) è un insieme di valori atomici, ovvero non ulteriormente scomponibili. 
Ogni dominio è caratterizzato da un nome e da un tipo di dato (es: \texttt{int, varchar, date}).
\begin{observation}[Il Tipo di Dato come Dominio Implicito]
Nella transizione dalla teoria all'ingegneria del software, il \textbf{tipo di dato} (es. \texttt{int}, \texttt{varchar}, \texttt{date}) definisce l'insieme universo di partenza, agendo come \textbf{dominio massimo implicito} dell'attributo. 

Eventuali restrizioni logiche (come un \textit{range} di valori validi per un voto universitario o un'età) si configurano come un \textbf{sottoinsieme} di tale dominio base. Vedremo in seguito come vengono implementate limitazioni specifiche. 
\end{observation}

\medskip

Uno \textbf{schema relazionale} (o schema di tabella) definisce la struttura logica dei dati. Esso è costituito da un nome $R$ e da un insieme di attributi $A_1, A_2, \dots, A_n$, a ciascuno dei quali è associato un dominio $D_i$, che denoteremo anche con $dom(A_i)$.

Lo schema viene denotato formalmente come:
$$R(A_1: D_1, A_2: D_2, \dots, A_n: D_n)$$

Lo schema relazionale corrisponde dunque esattamente al concetto di \textbf{Classe}. Esso definisce la struttura dei dati, i nomi dei campi e i rispettivi tipi, agendo come un ``blueprint'' statico.

Matematicamente, date le premesse precedenti, definiamo il dominio totale dello schema come il prodotto cartesiano dei domini dei singoli attributi:
$$D_{tot} = D_1 \times D_2 \times \dots \times D_n$$

\medskip

Una \textbf{relazione} $r$ sullo schema $R$ è un sottoinsieme finito del prodotto cartesiano dei domini degli attributi:
$$r \subseteq D_1 \times D_2 \times \dots \times D_n$$

Gli elementi che compongono questo sottoinsieme sono detti \textbf{tuple} (o n-uple). Una tupla è un vettore di valori $(v_1, v_2, \dots, v_n)$ tale che ogni $v_i \in D_i$.

\begin{tip}[]
La corrispondenza con OOP è la seguente:
\begin{itemize}
    \item Una \textbf{Tupla} corrisponde al singolo \textbf{Oggetto} (una istanza della classe contenente dati reali).
    \item La \textbf{Relazione} corrisponde all'insieme di tutti gli oggetti di quella classe \textbf{attualmente persistiti} nel database in un determinato istante temporale (lo stato corrente della tabella).
\end{itemize}
\end{tip}

\begin{example}[Formalizzazione di un'Anagrafica]
Definiamo lo schema relazionale per la classe \texttt{Studente}:
$$STUDENTE(\text{matricola}: \text{String}, \text{ nome}: \text{String}, \text{ eta}: \text{Integer})$$

Il prodotto cartesiano di questi tre domini genererebbe infinite combinazioni (comprese matricole inesistenti accoppiate a nomi casuali ed età assurde). 

La \textbf{relazione} $r(\text{STUDENTE})$ conterrà invece solo le tuple degli studenti reali inseriti nel sistema:
$$r = \{(\text{'654321'}, \text{'Giovanni'}, 22), (\text{'123456'}, \text{'Mario'}, 24)\}$$
\end{example}

\begin{observation}
    È significativo evidenziare che una relazione \textit{è un insieme}, dunque \textit{non ammette duplicati e non rappresenta una struttura ordinata}. 
\end{observation}

\begin{tip}[In sintesi:]
Nel passaggio dalla progettazione concettuale alla modellazione logica, i concetti del paradigma orientato agli oggetti (\textit{OOP}) e del Class Diagram si traducono come segue:

\begin{itemize}
    \item \textbf{Classi $\rightarrow$ Schemi Relazionali:} Le classi si traducono in schemi di relazione, identificati da un nome (corrispondente al nome della classe) e da un insieme di attributi.
    \item \textbf{Attributi $\rightarrow$ Domini:} A ciascun attributo dello schema è associato un \textbf{dominio}, ovvero l'insieme di tutti i possibili valori atomici che quell'attributo può assumere (equivalente al tipo di dato in programmazione).
    \item \textbf{Istanze $\rightarrow$ Relazione (Stato):} Una \textbf{relazione} è un sottoinsieme discreto e finito del dominio totale (il prodotto cartesiano di tutti i singoli domini). Essa rappresenta l'insieme delle tuple (istanze o oggetti) effettivamente memorizzate e persistite nel sistema in un dato momento temporale.
\end{itemize}
\end{tip}

% Consideriamo ora, nel class diagram, l'associazione studente-esame (uno a molti). A partire da una istanza di esame, per risalire allo studente che ha sostenuto quell'esame uso la chiave esterna matricola, da cui risalgo all'identificatore dello studente (matr, chiave primaria di STUDENTE).

% Da sottolineare quando si vanno rendere le associazioni molti a molti:

% sono problematiche perché il modello relazionale mi impone un valore unico per ciascun attributo, quindi non posso usare lo stesso metodo che uso per le associazioni uno a molti. 

% Il metodo per le associazioni molti a molti è del tutto generale ed in teoria si potrebbe usare per tutti i tipi di associazione; tuttavia questo metodo si usa solo dove strettamente necessario, perché ha un costo. 

\section{Le Chiavi}

Fino ad ora abbiamo analizzato come tradurre classi e attributi del Class Diagram nel modello relazionale, lasciando in sospeso la traduzione delle associazioni. Per poter modellare le relazioni tra tabelle, è prima necessario introdurre il concetto di chiave.

\medskip

Sia $R(A_1, A_2, \dots, A_n)$ uno schema di relazione. Un sottoinsieme di attributi $SK \subseteq \{A_1, A_2, \dots, A_n\}$ si definisce \textbf{superchiave} se i suoi valori identificano in modo univoco ogni singola tupla all'interno di qualsiasi istanza valida della relazione. 

\begin{example}[Superchiavi di un'Anagrafica]
Si consideri lo schema: 
$$\texttt{STUDENTE}({\text{CF}}, \text{ matricola}, \text{ nome}, \text{ cognome}, \text{ dataNascita})$$
I singoletti $\{\text{CF}\}$ e $\{\text{matricola}\}$ sono superchiavi. Di conseguenza, \textbf{qualsiasi sottoinsieme che li contenga} (es: $\{\text{CF}, \text{nome}\}$ o $\{\text{matricola}, \text{cognome}\}$) è a sua volta una superchiave.
\end{example}

\begin{observation}[Ipotesi di Esistenza]
Nella fase di progettazione concettuale si assume l'ipotesi che ogni istanza di una classe sia distinguibile dalle altre. Questa proprietà ci garantisce che \textbf{esiste sempre almeno una superchiave} per ogni schema relazionale: nella peggiore delle ipotesi (assenza di identificatori naturali), la superchiave coinciderà con l'insieme di tutti gli attributi dello schema.
\end{observation}

\medskip

Un sottoinsieme di attributi $K$ è una \textbf{chiave candidata} se è una \textbf{superchiave minimale}. Dire che una superchiave è \textit{minimale} significa che è \textbf{irriducibile}: eliminando un qualsiasi attributo da $K$, l'insieme rimanente perde la proprietà di univocità e cessa di essere una superchiave. 

Questo \textbf{non} significa che debba avere il numero minimo di attributi in assoluto (cardinalità minima). Se nello schema coesistono una superchiave irriducibile di cardinalità 1 e una di cardinalità 2, sono entrambe, con pari dignità, chiavi candidate.

È evidente che una superchiave composta da un unico attributo (un singoletto) è di fatto e necessariamente una chiave candidata.

\begin{example}[Chiave Candidata Composta]
Si consideri lo schema che mappa la classe degli esami universitari:
$$\texttt{ESAME}(\text{matricola}, \text{ corso}, \text{ voto}, \text{ data})$$
In questo scenario, il sottoinsieme $\{\text{matricola}, \text{corso}\}$ rappresenta una chiave candidata (di fatto, è l'unica chiave candidata in questo caso). Essa è minimale perché se eliminassimo \texttt{corso}, la sola \texttt{matricola} non garantirebbe l'univocità (uno studente può sostenere più esami); viceversa, eliminando \texttt{matricola}, il solo \texttt{corso} conterrebbe i voti di tutti gli studenti.
\end{example}

\medskip

Poiché in uno schema relazionale possono coesistere più chiavi candidate, il progettista deve sceglierne una da eleggere come identificatore ufficiale delle tuple. La chiave candidata scelta prende il nome di \textbf{Chiave Primaria} (\textit{Primary Key} - PK).

La selezione della chiave primaria è una decisione puramente ingegneristica e architetturale. Nello schema relazionale scritto, la chiave primaria si indica \textbf{sottolineando} gli attributi che la compongono.

\begin{example}[Scelta della PK]
Riprendendo lo schema \texttt{STUDENTE}, le chiavi candidate sono \texttt{CF} e \texttt{matricola}. Sceglieremo come chiave primaria la \texttt{matricola}:
$$\texttt{STUDENTE}(\text{CF}, \underline{\text{matricola}}, \text{nome}, \text{cognome})$$
\end{example}

\medskip

La \textbf{Chiave Esterna} (\textit{Foreign Key} - FK) è lo strumento fondamentale del modello relazionale per tradurre le associazioni.

Siano $R_1$ e $R_2$ due schemi relazionali (non necessariamente distinti). Un insieme di attributi $FK$ di $R_1$ è una chiave esterna che punta a $R_2$ se:
\begin{enumerate}
    \item Gli attributi in $FK$ sono in corrispondenza biunivoca con gli attributi che compongono la \textbf{chiave primaria} ($PK$) di $R_2$.
    \item I domini degli attributi corrispondenti sono identici ($dom(FK_i) = dom(PK_i)$).
\end{enumerate}

\medskip

Il vincolo impone che per ogni tupla $t_1$ nell'istanza corrente di $R_1$, il valore di $t_1[FK]$ deve essere nullo (\texttt{NULL}), oppure coincidere esattamente con il valore $t_2[PK]$ di \textbf{una tupla} $t_2$ \textbf{attualmente esistente} nell'istanza di $R_2$.

Nello schema testuale, la presenza di una chiave esterna si denota convenzionalmente con una \textbf{doppia sottolineatura}.

\section{Traduzione delle Associazioni}

\subsection{Associazioni 1 a 1}

Un'associazione \textbf{1 a 1} tra due classi $A$ e $B$ indica che a un'istanza di $A$ corrisponde al massimo un'istanza di $B$, e viceversa. 
Nel modello relazionale, questo legame si traduce introducendo una \textbf{Chiave Esterna (FK)} in uno dei due schemi che punta alla Chiave Primaria (PK) dell'altro.

\begin{example}[Associazione esattamente 1 a 1]
Si consideri un'associazione in cui la partecipazione è rigida e obbligatoria per entrambe le classi (1..1 $\leftrightarrow$ 1..1), ad esempio tra l'entità \texttt{STUDENTE(matricola, nome, cognome)} e il suo \texttt{FASCICOLO(codice, dataCreazione)}. Ogni studente possiede esattamente un fascicolo e ogni fascicolo appartiene esattamente a uno studente.
Traduciamo il modello aggiungendo la chiave esterna nello schema del fascicolo:
\begin{itemize}
    \item $\texttt{STUDENTE}(\underline{\text{matricola}}, \text{ nome}, \text{ cognome})$
    \item $\texttt{FASCICOLO}(\underline{\text{codice}}, \text{ dataCreazione}, \underline{\underline{\text{ matricolaStudente}}})$
\end{itemize}
\end{example}

Affinché la relazione rimanga rigorosamente 1 a 1, non è sufficiente definire la chiave esterna. Se ci limitassimo a questo, la tabella che ospita la FK potrebbe avere più tuple che puntano alla stessa PK della tabella bersaglio, trasformando di fatto la relazione in una 1 a N.

Riprendendo l'esempio precedente, se l'attributo \texttt{studente} nello schema \texttt{FASCICOLO} fosse definito unicamente come chiave esterna, nulla impedirebbe al database di registrare due tuple distinte di fascicoli (es. codice \texttt{'F01'} e \texttt{'F02'}) aventi entrambe il valore \texttt{studente = '12345'}. Questo corromperebbe la semantica concettuale, asserendo che la matricola 12345 possiede \textit{molti} fascicoli.

Per garantire la biunivocità dell'associazione, è obbligatorio imporre un \textbf{vincolo di univocità} sugli attributi che compongono la chiave esterna.
Matematicamente, se lo schema $R_1$ contiene la chiave esterna $FK$ verso $R_2$, deve valere che per ogni coppia di tuple distinte $t_x, t_y \in r(R_1)$:
$$t_x[FK] \neq t_y[FK] \quad \text{(escludendo i valori NULL)}$$

\begin{dettagli}[Perché usare una sola Chiave Esterna?]
Vista la forte simmetria di un'associazione 1 a 1, potrebbe sorgere spontaneo il dubbio: perché non posizionare una chiave esterna in \textit{entrambi} gli schemi (es. aggiungere \texttt{fascicolo} dentro \texttt{STUDENTE} e \texttt{studente} dentro \texttt{FASCICOLO})? 
La teoria relazionale esclude questa pratica per tre ragioni fondamentali:
\begin{enumerate}
    \item \textbf{Ridondanza e Rischio di Anomalie:} Il legame logico verrebbe salvato due volte. Se per errore un sistema aggiornasse un lato senza allineare l'altro (es. lo Studente punta al Fascicolo A, ma il Fascicolo A punta a un altro Studente), il database entrerebbe in uno stato incoerente.
    \item \textbf{Bidirezionalità Intrinseca:} Come già stabilito, la presenza di una singola chiave esterna è matematicamente sufficiente per stabilire un'uguaglianza logica tra le tuple. Il motore del database può navigare il legame in entrambe le direzioni senza bisogno di puntatori doppi.
    \item \textbf{Dipendenza Circolare (Deadlock di Inserimento):} Poiché in questo esempio la partecipazione è obbligatoria da ambo i lati, le rispettive chiavi esterne dovrebbero avere un vincolo \texttt{NOT NULL}. Tuttavia, questo creerebbe un paradosso logico: non si potrebbe inserire uno studente nel database perché manca il suo fascicolo, e non si potrebbe inserire il fascicolo perché manca lo studente a cui associarlo. Inserire la chiave esterna in un solo schema spezza questo circolo vizioso.
\end{enumerate}
\end{dettagli}

\begin{observation}
    In alcuni casi estremi di forte coesione, le due tabelle possono persino essere fuse in un unico schema relazionale. Approfondiremo questo aspetto nel prossimo paragrafo. 
\end{observation}

Nell'esempio precedente si aveva una situazione di relazione esattamente 1 a 1 e la scelta dello schema in cui inserire la chiave esterna era indifferente. In altri casi, la scelta di quale delle due tabelle debba ospitare la FK dipende dalle \textbf{cardinalità minime} (l'opzionalità o obbligatorietà dell'associazione): se la classe $A$ partecipa in modo opzionale (0..1) e la classe $B$ in modo obbligatorio (1..1), la chiave esterna \textbf{deve essere inserita nello schema $B$}. In questo modo, la $FK$ in $B$ potrà essere dichiarata \texttt{NOT NULL} (vincolo di esistenza garantito). Se la mettessimo in $A$, saremmo costretti ad ammettere valori \texttt{NULL} per le istanze di $A$ che non partecipano all'associazione, sprecando spazio e riducendo l'eleganza relazionale.

\begin{example}[Direzione di un Dipartimento]
Si consideri l'associazione 1 a 1 tra \texttt{Professore} (0..1) e \texttt{Dipartimento} (1..1): un professore può dirigere al massimo un dipartimento, ma un dipartimento deve avere esattamente un direttore.

Seguendo la regola, posizioniamo la chiave esterna nel \texttt{Dipartimento}:
\begin{itemize}
    \item $\texttt{PROFESSORE}(\underline{\text{codiceDocente}}, \text{nome}, \text{cognome})$
    \item $\texttt{DIPARTIMENTO}(\underline{\text{codice}}, \text{nome}, \underline{\underline{\text{codiceDirettore}}})$
\end{itemize}

\medskip

Vincoli su \texttt{DIPARTIMENTO.codiceDirettore}:
\begin{itemize}
    \item \textbf{Foreign Key:} punta a \texttt{PROFESSORE.codiceDocente}
    \item \textbf{UNIQUE:} per garantire che lo stesso professore non diriga due dipartimenti (vincolo 1:1).
    \item \textbf{NOT NULL:} per garantire che ogni dipartimento abbia un direttore (cardinalità minima 1).
\end{itemize}
\end{example}

\begin{dettagli}[Controllo Applicativo vs Vincolo Strutturale (UNIQUE)]
In alcuni contesti didattici o in framework di sviluppo specifici, è possibile omettere la dichiarazione esplicita del vincolo \texttt{UNIQUE} per le associazioni 1 a 1. 

In questi scenari, si assume un'ipotesi di \textbf{coerenza applicativa}: si dà per scontato che sia la logica del software sovrastante (es. i controlli nel codice backend) a impedire l'inserimento di tuple multiple, in accordo con la natura del dominio (``non può mai capitare che l'entità sia associata a molti'').

Tuttavia, dal punto di vista della solidità ingegneristica e dell'integrità fisica dei dati (\textit{defensive programming}), l'omissione del vincolo \texttt{UNIQUE} rende il database strutturalmente identico a un'associazione 1 a Molti. Fissare il vincolo a livello di DBMS rimane la pratica raccomandata per garantire che i dati non vengano mai corrotti, indipendentemente da eventuali bug del software applicativo.
\end{dettagli}

\subsection{Associazioni 1 a Molti (1:N)}

Nel modello relazionale, questa associazione si traduce secondo una regola fissa e inderogabile, dettata dal vincolo di atomicità degli attributi (assenza di attributi multivalore): la \textbf{Chiave Esterna (FK) deve essere sempre inserita nello schema dal lato ``Molti'' ($B$)}, e punterà alla Chiave Primaria (PK) dello schema dal lato ``Uno'' ($A$).

A differenza dell'associazione 1 a 1, in questo caso \textbf{non} si deve applicare alcun vincolo di univocità (\texttt{UNIQUE}) sulla chiave esterna. È infatti intrinseco nella natura dell'associazione 1:N che più tuple della relazione $B$ possano possedere lo stesso valore per l'attributo FK, puntando così alla medesima istanza di $A$.

\begin{example}[Impiegati in un Dipartimento]
Si consideri l'associazione tra \texttt{Dipartimento} (1) e \texttt{Impiegato} (N). Un dipartimento ospita molti impiegati, ma un impiegato afferisce a un solo dipartimento.
La chiave esterna viene posizionata nello schema \texttt{IMPIEGATO}:

\begin{itemize}
    \item $\texttt{DIPARTIMENTO}(\underline{\text{codice}}, \text{nome}, \text{sede})$
    \item $\texttt{IMPIEGATO}(\underline{\text{matricola}}, \text{nome}, \text{cognome}, \underline{\underline{\text{codiceDipartimento}}})$
\end{itemize}
\end{example}

\begin{observation}[Monodirezionalità Apparente]
Anche in questo caso, la posizione fisica della chiave esterna genera una \textit{monodirezionalità apparente}. Guardando lo schema, sembra che un impiegato "conosca" il proprio dipartimento, ma che il dipartimento sia ignaro dei propri impiegati (poiché non ha attributi che li referenziano). 

Come già chiarito, questa è solo un'illusione strutturale: l'informazione è conservata intatta dalla logica dei valori. Il database è sempre in grado di ricavare l'elenco degli impiegati di uno specifico dipartimento, semplicemente interrogando la tabella \texttt{IMPIEGATO} e filtrando le tuple che possiedono quel determinato \texttt{codiceDipartimento}.
\end{observation}

\subsection{Associazioni Molti a Molti (N:N)}

Nei casi precedenti (1:1 e 1:N), la traduzione avveniva inserendo una chiave esterna direttamente in uno dei due schemi coinvolti. 
Se provassimo ad applicare questa logica a un'associazione N:N, violeremmo il vincolo fondamentale del modello relazionale:
\begin{itemize}
    \item Inserire la chiave esterna di \texttt{CORSO} nello schema \texttt{STUDENTE} costringerebbe una tupla studente a contenere un attributo \textbf{multivalore} (una lista di codici di corsi).
    \item Viceversa, inserire la chiave esterna di \texttt{STUDENTE} in \texttt{CORSO} creerebbe una lista di matricole dentro la tupla del corso.
\end{itemize}
Poiché il modello relazionale ammette esclusivamente valori atomici, l'uso di chiavi esterne dirette è matematicamente impossibile.

\medskip

Per tradurre un'associazione Molti a Molti, la teoria relazionale impone la creazione di un \textbf{terzo schema di relazione autonomo}, comunemente chiamato \textit{tabella di giunzione}, \textit{tabella ponte} o \textit{relazione associativa}.

Questa nuova tabella conterrà esclusivamente (almeno nella sua forma base) le chiavi esterne che puntano alle chiavi primarie dei due schemi originari. La \textbf{Chiave Primaria} di questa nuova tabella sarà composta dall'unione delle due chiavi esterne (chiave primaria composta).

\begin{example}[Studenti e Corsi]
Si consideri l'associazione N:N tra \texttt{STUDENTE} e \texttt{CORSO}.
Creiamo la tabella ponte \texttt{FREQUENZA}:

\begin{itemize}
    \item $\texttt{STUDENTE}(\underline{\text{matricola}}, \text{nome})$
    \item $\texttt{CORSO}(\underline{\text{codice}}, \text{titolo}, \text{cfu})$
    \item $\texttt{FREQUENZA}(\underline{\underline{\text{matricolaStudente}}}, \underline{\underline{\text{codiceCorso}}})$
\end{itemize}

\medskip

Vincoli sulla tabella \texttt{FREQUENZA}:
\begin{itemize}
    \item \textbf{Primary Key:} È composta dalla coppia (\texttt{matricolaStudente}, \texttt{codiceCorso}). Questo garantisce che uno studente non possa essere iscritto due volte allo stesso identico corso.
    \item \textbf{Foreign Key 1:} \texttt{matricolaStudente} punta a \texttt{STUDENTE.matricola}.
    \item \textbf{Foreign Key 2:} \texttt{codiceCorso} punta a \texttt{CORSO.codice}.
\end{itemize}
\end{example}

\begin{observation}[Delegazione della Complessità]
Dal punto di vista della progettazione (struttura statica), il pattern della tabella di giunzione rende la risoluzione delle associazioni N:N un processo banale e standardizzato.
La reale complessità viene delegata alla fase operativa (interrogazione dei dati). Per estrarre, ad esempio, i nomi degli studenti e i titoli dei corsi da loro frequentati, il DBMS dovrà eseguire computazionalmente onerosi \textbf{JOIN multipli}, attraversando la tabella ponte per ricongiungere i dati disaccoppiati.

È chiaro che questo metodo per rendere le associazioni è del tutto generale e, in teoria, potrebbe essere usato per tutti i tipi di associazione; tuttavia, si tratta ovviamente di una cattiva pratica in virtù del costo computazionale che è stato appena discusso. 
\end{observation}

\subsubsection*{Classi di Associazione}

Nel modellamento concettuale tramite Class Diagram UML, capita frequentemente che un'associazione molti a molti possieda a sua volta degli attributi (formalizzati tramite una \textbf{Classe di Associazione}). Un esempio tipico è la relazione di superamento di un esame tra uno studente e un corso, la quale porta con sé attributi come il \textit{voto}, la \textit{data} e la \textit{lode}.

La traduzione di questo costrutto nel modello relazionale segue fedelmente la regola dello schema ponte vista in precedenza. Gli attributi propri dell'associazione vengono semplicemente inseriti come \textbf{colonne aggiuntive all'interno della tabella di giunzione}.

\begin{example}[Mappatura della Classe di Associazione Esame]
Estendiamo l'esempio precedente trasformando lo schema ponte \texttt{FREQUENZA} in uno schema \texttt{ESAME} che ospiti anche i dati del superamento del corso:

\begin{itemize}
    \item $\texttt{STUDENTE}(\underline{\text{matricola}}, \text{nome}, \text{cognome})$
    \item $\texttt{CORSO}(\underline{\text{codice}}, \text{titolo}, \text{cfu})$
    \item $\texttt{ESAME}(\underline{\underline{\text{matricolaStudente}}}, \underline{\underline{\text{codiceCorso}}}, \text{voto}, \text{data}, \text{lode})$
\end{itemize}

\medskip

Nello schema \texttt{ESAME}, la chiave primaria rimane la coppia composta (matricolaStudente, codiceCorso), il che modella il vincolo per cui uno studente può registrare al massimo un voto per un determinato corso. Gli attributi \texttt{voto}, \texttt{data} e \texttt{lode} diventano attributi non chiave della tabella di giunzione.
\end{example}

\subsubsection{Associazioni n-arie}

La logica della tabella di giunzione, introdotta per le associazioni binarie Molti a Molti, scala in modo naturale e identico per le \textbf{associazioni n-arie} (con $n > 2$). 

Quando nel diagramma concettuale un'associazione coinvolge tre o più classi (ad esempio un'associazione ternaria), la traduzione nel modello relazionale richiede la creazione di un unico \textbf{schema ponte} dedicato. Questo schema assolverà il compito di legare tutte le entità simultaneamente e conterrà:
\begin{itemize}
    \item Una \textbf{Chiave Esterna (FK)} per ciascuna delle $n$ entità partecipanti, che punta alla rispettiva chiave primaria.
    \item Una \textbf{Chiave Primaria (PK) composta} dall'unione di tutte le chiavi esterne coinvolte (salvo casi particolari in cui vincoli di cardinalità massima pari a 1 su specifici rami consentano di escludere alcune FK dalla chiave primaria).
    \item Gli eventuali attributi dell'associazione (o della Classe di Associazione), inseriti semplicemente come colonne non chiave.
\end{itemize}

\begin{example}[Associazione Ternaria di Fornitura]
Si consideri un'associazione ternaria che modella la fornitura di materiali aziendali: un determinato \texttt{Fornitore} vende uno specifico \texttt{Prodotto} destinato a un particolare \texttt{Progetto}, in una certa \texttt{quantita}.

La traduzione genera tre schemi base e un quarto schema ponte per l'associazione ternaria:
\begin{itemize}
    \item $\texttt{FORNITORE}(\underline{\text{partitaIva}}, \text{nomeAzienda})$
    \item $\texttt{PRODOTTO}(\underline{\text{codiceProdotto}}, \text{descrizione})$
    \item $\texttt{PROGETTO}(\underline{\text{idProgetto}}, \text{budget})$
    \item $\texttt{FORNITURA}(\underline{\underline{\text{partitaIva}}}, \underline{\underline{\text{codiceProdotto}}}, \underline{\underline{\text{idProgetto}}}, \text{quantita})$
\end{itemize}

Nello schema \texttt{FORNITURA}, la chiave primaria è la tripla di attributi (partitaIva, codiceProdotto, idProgetto). L'attributo \texttt{quantita} è un attributo dell'associazione che dipende funzionalmente dall'intera tripla, in quanto esprime quanti pezzi di quel prodotto specifico sono stati comprati da quel fornitore per quel preciso progetto.
\end{example}

\section{Ristrutturazione dello Schema Concettuale}

Il Class Diagram realizzato durante la progettazione concettuale ha l'unico scopo di modellare la semantica del dominio applicativo, ignorando i limiti tecnologici. Tuttavia, avvicinandosi alla progettazione logica, l'adozione dello specifico modello relazionale fa emergere vincoli strutturali e di performance che rendono il diagramma originario inefficiente o inapplicabile. 

Prima di procedere alla traduzione in tabelle, è quindi necessaria una fase di \textbf{Ristrutturazione del Class Diagram}, durante la quale il modello viene riadattato e "sporcato" con dettagli implementativi. Questa fase si articola in diverse analisi sequenziali.

\subsection{Analisi delle Chiavi: Chiavi Naturali vs Surrogate}

La prima operazione consiste nel verificare gli identificatori (chiavi primarie) scelti per le classi. Nel modello concettuale, è prassi utilizzare come identificatori degli attributi già esistenti nel dominio della realtà, noti come \textbf{Chiavi Naturali}. 
Tuttavia, capita spesso che per identificare univocamente un'istanza nel mondo reale sia necessario unire molti attributi, creando una \textbf{chiave primaria composta} molto ingombrante.

L'uso di chiavi primarie formate da tre, quattro o più attributi genera due enormi problemi nel database relazionale:
\begin{enumerate}
    \item \textbf{Propagazione delle Foreign Key:} Poiché le associazioni si traducono copiando la chiave primaria nella tabella collegata, una chiave di 4 attributi costringerà tutte le tabelle figlie ad avere 4 colonne aggiuntive solo per mantenere il legame. Questo allarga a dismisura le tabelle e spreca memoria.
    \item \textbf{Costo Computazionale:} Gli indici sulle chiavi composite sono lenti da scorrere e pesanti da aggiornare, rendendo i \texttt{JOIN} estremamente inefficienti.
\end{enumerate}

\medskip

Per risolvere questo problema, si modifica il Class Diagram introducendo un nuovo attributo artificiale, privo di alcun significato nel dominio di business, chiamato \textbf{Chiave Surrogata} (o \textit{Sintetica}).
Generalmente si tratta di un numero intero auto-incrementante (es. \texttt{id}, \texttt{codice}) o di un UUID alfanumerico. Questo attributo viene eletto a nuova Chiave Primaria, sostituendo la chiave naturale composta.

\begin{example}[Introduzione di una Chiave Surrogata]
Si consideri la classe \texttt{Stanza} di un polo universitario. Nel mondo reale, una stanza è identificata univocamente dalla combinazione di tre attributi: \texttt{edificio}, \texttt{piano} e \texttt{numeroStanza}.
$$\texttt{STANZA}(\underline{\text{edificio}}, \underline{\text{piano}}, \underline{\text{numeroStanza}}, \text{capienza})$$
Se un corso dovesse essere assegnato a una stanza, lo schema \texttt{CORSO} dovrebbe ereditare tutte e tre le colonne come Foreign Key.

Ristrutturiamo lo schema introducendo una chiave surrogata \texttt{idStanza}:
$$\texttt{STANZA}(\underline{\text{idStanza}}, \text{edificio}, \text{piano}, \text{numeroStanza}, \text{capienza})$$
Ora, le tabelle collegate (come \texttt{CORSO}) dovranno importare solamente la colonna \texttt{idStanza} (un semplice intero) per stabilire l'associazione, ottimizzando drasticamente le prestazioni del database.
\end{example}

\begin{warningbox}[Attenzione alla perdita di integrità]
Quando si sostituisce una chiave naturale con una chiave surrogata, il database userà il nuovo \texttt{id} per distinguere le righe. Questo crea un potenziale rischio: il database permetterà di inserire due righe con \texttt{id} diversi ma con gli \textbf{stessi identici dati naturali} (es. due stanze diverse assegnate allo stesso edificio, stesso piano e stesso numero). 
Per evitare dati duplicati, è \textbf{obbligatorio} imporre un vincolo di univocità (\texttt{UNIQUE}) sull'insieme degli attributi che formavano l'originaria chiave naturale composita.
\end{warningbox}

\subsection{Analisi degli Attributi Derivati e delle Ridondanze}

Durante la progettazione concettuale, è comune inserire attributi il cui valore non è indipendente, ma può essere ricavato matematicamente o logicamente da altri dati presenti nel sistema. Nel Class Diagram UML, questa caratteristica viene esplicitata anteponendo una barra (slash) al nome dell'attributo (es. \texttt{/eta}).

Nel passaggio al Class Diagram revisionato (e successivamente allo schema logico), il progettista deve compiere una scelta architetturale precisa: l'attributo derivato verrà \textbf{effettivamente memorizzato} (storicizzato) nel database o verrà eliminato dallo schema e \textbf{calcolato a runtime} ogni volta che viene richiesto? Se si decide di non memorizzarlo, l'attributo scompare del tutto dallo schema revisionato.

\begin{example}[Età vs Data di Nascita]
Si consideri la classe \texttt{Persona} dotata degli attributi \texttt{dataDiNascita} ed \texttt{/eta}. L'età è banalmente derivabile per differenza tra la data odierna e la data di nascita.
\end{example}

La scelta tra memorizzazione e calcolo on-the-fly comporta un compromesso tra prestazioni di lettura, spazio occupato e coerenza dei dati.

\medskip

\textbf{Scenario A: Memorizzazione (Storicizzazione del dato)}
\begin{itemize}
    \item \textbf{Vantaggi:} Tempi di risposta immediati durante le interrogazioni (query). Il DBMS si limita a leggere il dato senza sprecare cicli di CPU per eseguire calcoli, aspetto cruciale se il dato è richiesto frequentissimamente.
    \item \textbf{Svantaggi:} 
    \begin{itemize}
        \item Spreco di spazio su disco (memoria).
        \item Rischio di \textit{inconsistenza}: se la \texttt{dataDiNascita} viene corretta ma ci si dimentica di aggiornare l'\texttt{eta}, il database conterrà informazioni contraddittorie.
        \item Oneri di aggiornamento: dati altamente volatili (come l'età, che cambia ogni anno) richiedono procedure o \textit{trigger} che aggiornino continuamente il database, generando un carico di scrittura pesante.
    \end{itemize}
\end{itemize}

\medskip

\textbf{Scenario B: Calcolo a Runtime (Nessuna memorizzazione)}
\begin{itemize}
    \item \textbf{Vantaggi:} Nessun rischio di inconsistenza (il calcolo parte sempre dal "dato sorgente" aggiornato), risparmio di memoria e zero manutenzione per l'aggiornamento.
    \item \textbf{Svantaggi:} Costo computazionale pagato al momento dell'interrogazione. Se il calcolo è complesso e viene richiesto da milioni di utenti, il server potrebbe rallentare drasticamente.
\end{itemize}

\begin{tip}[Principi Generali per gli Attributi Derivati]
Nell'ingegneria dei dati, la scelta si basa sull'analisi del carico di lavoro previsto e sulla volatilità del dato:
\begin{enumerate}
    \item \textbf{Dati volatili e calcoli leggeri (Preferire il Calcolo):} Se il dato derivato cambia spesso nel tempo e l'operazione matematica è banale, l'attributo si elimina e si calcola al volo.
    \item \textbf{Dati stabili e calcoli pesanti (Preferire la Memorizzazione):} Se il dato originario, una volta inserito, non cambia più (es. il totale di una \texttt{Fattura} emessa e consolidata, ricavato dalla somma delle sue \texttt{Righe}), conviene storicizzare il totale. Il costo di aggiornamento è zero (la fattura non cambierà) e si evitano pesanti \texttt{JOIN} e somme ad ogni lettura.
\end{enumerate}
\end{tip}

\subsubsection*{Analisi delle Associazioni Ridondanti}

Il concetto di ridondanza si estende oltre i singoli attributi e coinvolge l'intera topologia del Class Diagram. In uno schema complesso, potremmo accorgerci dell'esistenza di \textbf{associazioni derivabili} per transitività. 

Ad esempio, se un \texttt{Impiegato} è associato a un \texttt{Dipartimento}, e un \texttt{Dipartimento} è associato a una \texttt{Città}, un'eventuale associazione diretta tra \texttt{Impiegato} e \texttt{Città} è da considerarsi un ciclo ridondante, poiché la città dell'impiegato è già ricavabile navigando il percorso \texttt{Impiegato $\rightarrow$ Dipartimento $\rightarrow$ Città}.

Come per gli attributi, il progettista deve decidere se mantenere l'associazione diretta (tramite l'aggiunta di una chiave esterna ridondante) o se farla collassare e affidarsi alla navigazione del grafo.

\begin{tip}[Principi Generali per le Associazioni Ridondanti]
La gestione dei cicli si basa sul compromesso tra prestazioni di \textit{routing} e integrità strutturale:
\begin{enumerate}
    \item \textbf{Evitare anomalie strutturali (Preferire l'eliminazione):} La regola generale impone di rimuovere le associazioni ridondanti. Mantenere un legame diretto crea il rischio di inconsistenze: potremmo, per errore, assegnare l'Impiegato al Dipartimento di Milano, ma associarlo direttamente alla Città di Roma. 
    \item \textbf{Ottimizzazione estrema (Mantenere il legame diretto):} Si deroga alla regola precedente solo in sistemi ad altissimo traffico (\textit{read-heavy}) in cui il percorso di navigazione è molto lungo (es. richiede di attraversare 4 o 5 tabelle con relative operazioni di \texttt{JOIN}). In quel caso, si introduce l'associazione ridondante come scorciatoia prestazionale, assumendosi l'onere (spesso tramite \textit{stored procedures}) di mantenerla sincronizzata con il percorso principale.
\end{enumerate}
\end{tip}

\subsection{Analisi delle Associazioni Uno a Uno: Fusione vs Separazione}

Durante la progettazione concettuale, due entità legate da un'associazione rigorosamente biunivoca (1..1 $\leftrightarrow$ 1..1) mantengono identità separate per ragioni semantiche. Poiché i loro cicli di vita coincidono (nascono e muoiono insieme), nella fase di ristrutturazione logica si presenta l'opportunità di \textbf{fondere le due classi in un unico schema relazionale}, consolidando i loro attributi ed eliminando la necessità di una chiave esterna.

Questa decisione architetturale comporta un delicato trade-off tra la velocità di lettura (assenza di \texttt{JOIN}) e l'occupazione di memoria (pesantezza della tabella), oltre a sollevare questioni di sicurezza dei dati.

\begin{example}[Fusione vs Separazione: Paziente e Cartella Clinica]
Si consideri l'associazione obbligatoria (1..1 $\leftrightarrow$ 1..1) tra un \texttt{Paziente} e la sua \texttt{CartellaClinica}. Ogni paziente ha una sola cartella clinica, e ogni cartella appartiene a un solo paziente.

\textbf{Scenario A: Fusione (Unico Schema)}
Se fondessimo le due entità, otterremmo un'unica tabella massiccia:
$$\texttt{PAZIENTE}(\underline{\text{codiceFiscale}}, \text{nome}, \text{telefono}, \text{gruppoSanguigno}, \text{anamnesiCompleta})$$
\begin{itemize}
    \item \textit{Vantaggio:} Nessun \texttt{JOIN} necessario per leggere i dati completi.
    \item \textit{Svantaggio:} L'attributo \texttt{anamnesiCompleta} potrebbe essere un testo lunghissimo (BLOB o TEXT). Ogni volta che un impiegato dell'accettazione interroga il database solo per cercare il numero di \texttt{telefono} del paziente, il motore del database è costretto a caricare in memoria dal disco fisso anche i megabyte di testo dell'anamnesi, sprecando enormi risorse di I/O. Inoltre, si espongono dati sensibili a personale non medico.
\end{itemize}

\medskip

\textbf{Scenario B: Separazione (Due Schemi)}
Mantenendo gli schemi separati, applichiamo un \textit{partizionamento verticale} logico:
\begin{itemize}
    \item $\texttt{PAZIENTE}(\underline{\text{codiceFiscale}}, \text{nome}, \text{telefono})$
    \item $\texttt{CARTELLA\_CLINICA}(\underline{\text{codiceCartella}}, \text{gruppoSanguigno}, \text{anamnesi}, \underline{\underline{\text{cfPaziente}}})$
\end{itemize}

\medskip

\textit{Vantaggi:} 
\begin{enumerate}
    \item \textbf{Performance:} La tabella \texttt{PAZIENTE} è "leggera" e stretta. Le query anagrafiche quotidiane saranno fulminee e il database potrà tenere l'intera tabella nella cache RAM.
    \item \textbf{Sicurezza:} Si possono impostare permessi di accesso (GRANT) a livello di tabella, permettendo alla segreteria di leggere \texttt{PAZIENTE} e impedendo l'accesso a \texttt{CARTELLA\_CLINICA}, riservato ai soli medici.
\end{enumerate}

\medskip

\textit{Svantaggio:} Quando un medico avrà bisogno del nome del paziente e della sua anamnesi contemporaneamente, il sistema dovrà pagare il costo computazionale di un \texttt{JOIN}.
\end{example}

\begin{tip}[La Regola Pratica]
La fusione in un unico schema è raccomandata quando gli attributi della seconda classe sono pochi, di piccole dimensioni (es. interi, date, stringhe brevi) e vengono interrogati quasi sempre insieme a quelli della classe principale (alta coesione). La separazione è invece mandatoria in presenza di attributi "pesanti" (es. file, immagini, testi lunghi) o sottoposti a stringenti policy di privacy.
\end{tip}

\subsection{Analisi degli Attributi Strutturati}

Nel diagramma delle classi originario, è frequente l'utilizzo di \textbf{attributi strutturati} (o \textit{composti}), ovvero attributi che raggruppano al loro interno campi di natura eterogenea. L'esempio classico è l'attributo \texttt{residenza} assegnato a una classe \texttt{Persona}, il quale è concettualmente composto da \texttt{via}, \texttt{civico}, \texttt{cap} e \texttt{citta}.

Il Modello Relazionale, tuttavia, impone il rispetto della \textbf{Prima Forma Normale (1NF)}, la quale stabilisce che il dominio di ogni attributo debba essere \textbf{atomico} (indivisibile). I database relazionali puri non supportano l'inserimento di "sotto-strutture" o "oggetti" all'interno di una singola cella. 

Per risolvere questa incompatibilità durante la fase di ristrutturazione, il progettista deve eliminare l'attributo strutturato scegliendo una delle seguenti tre strategie:

\subsubsection*{1. Appiattimento (Flattening) degli attributi}
L'attributo strutturato viene rimosso e sostituito dall'esplicitazione di tutti i suoi sotto-campi direttamente come colonne della classe principale.
\begin{itemize}
    \item \textbf{Schema:} $\texttt{PERSONA}(\underline{\text{cf}}, \text{nome}, \text{via}, \text{civico}, \text{cap}, \text{citta})$
    \item \textbf{Vantaggi:} Non richiede \texttt{JOIN}. Permette di effettuare ricerche efficienti sui singoli sotto-campi (es. "trova tutte le persone che abitano a Roma").
    \item \textbf{Svantaggi:} "Allarga" la tabella principale, rischiando di riempirla di valori \texttt{NULL} qualora l'attributo strutturato originario fosse opzionale (0..1).
\end{itemize}

\subsubsection*{2. Accorpamento in formato Stringa (Concatenazione)}
I sotto-campi vengono fusi in un'unica stringa testuale, inserita in un solo attributo atomico.
\begin{itemize}
    \item \textbf{Schema:} $\texttt{PERSONA}(\underline{\text{cf}}, \text{nome}, \text{indirizzoCompleto})$
    \item \textbf{Vantaggi:} Estrema semplicità e compattezza dello schema.
    \item \textbf{Svantaggi:} Rende difficilissime o ineseguibili le query analitiche (es. contare le persone per CAP richiederebbe complesse e lente operazioni di \textit{parsing} testuale tramite \texttt{LIKE} o \textit{Regex}).
\end{itemize}

\subsubsection*{3. Creazione di un'Entità Autonoma}
Si estrapola l'attributo strutturato promuovendolo a nuova \textbf{Classe autonoma}, collegata a quella originale tramite un'associazione.
\begin{itemize}
    \item \textbf{Schema:} $\texttt{PERSONA}(\dots, \underline{\underline{\text{idResidenza}}}) \rightarrow \texttt{RESIDENZA}(\underline{\text{id}}, \text{via}, \text{civico}, \text{citta})$
    \item \textbf{Vantaggi:} Massima pulizia formale. Permette di riutilizzare l'entità (es. se cinque membri di una famiglia condividono la stessa residenza, il dato non viene duplicato 5 volte, ma inserito una volta sola e linkato).
    \item \textbf{Svantaggi:} Introduce un costo computazionale (\texttt{JOIN}) per ogni lettura completa e rende più complessa la fase di inserimento dei dati.
\end{itemize}

\begin{tip}[Principi Generali di Scelta]
La decisione tra queste tre soluzioni è dettata dall'utilizzo pratico che l'applicazione farà del dato:
\begin{enumerate}
    \item \textbf{Se il dato è un blocco opaco (Preferire l'Accorpamento testuale):} Se l'indirizzo serve esclusivamente per essere stampato su un'etichetta di spedizione e al sistema non interesserà mai fare statistiche per città o per CAP, una stringa formattata è la scelta più leggera e sensata.
    \item \textbf{Se i campi hanno dignità di interrogazione (Preferire l'Appiattimento):} Se il software prevede filtri di ricerca (es. un menù a tendina per filtrare gli utenti per Provincia), i campi devono essere separati per poter essere indicizzati singolarmente.
    \item \textbf{Se l'attributo ha vita propria o cardinalità multipla (Preferire l'Entità Autonoma):} Se l'indirizzo deve essere storicizzato (tenere traccia dei vecchi indirizzi), se una persona ha molteplici indirizzi (casa, lavoro, fatturazione), o se lo stesso indirizzo raggruppa logicamente più utenti (come in un'azienda), l'attributo strutturato \textbf{deve} diventare una tabella separata.
\end{enumerate}
\end{tip}

\subsection{Analisi degli Attributi a Valore Multiplo}

Nel Class Diagram UML è perfettamente lecito assegnare a un'entità un attributo con cardinalità multipla, ovvero una lista o un array di valori (es. \texttt{telefoni [1..*]} per la classe \texttt{Persona}). 
Tuttavia, il modello relazionale ammette unicamente \textbf{valori singoli e scalari} per ogni cella. L'inserimento di vettori o array non è nativamente supportato dalle regole dell'algebra relazionale.

Il tentativo di risolvere il problema "appiattendo" l'attributo in una serie di colonne fisse (\texttt{telefono1}, \texttt{telefono2}, \dots) rappresenta una pessima pratica ingegneristica: spreca enormi quantità di memoria (generando celle \texttt{NULL} per gli utenti con un solo recapito), limita la scalabilità (non permette l'inserimento di un numero di recapiti superiore alle colonne predisposte) e rende le query di ricerca estremamente inefficienti.

Per normalizzare lo schema, il progettista deve eliminare l'attributo multivalore ricorrendo a una delle seguenti soluzioni:

\subsubsection*{1. Accorpamento testuale (Formattazione Custom)}
Tutti i valori dell'array vengono concatenati e salvati in un'unica stringa di testo, separati da un carattere speciale (es. virgola o punto e virgola).
\begin{itemize}
    \item \textbf{Schema:} $\texttt{PERSONA}(\underline{\text{cf}}, \text{nome}, \text{listaTelefoni})$ 
    \item \textit{Valore di esempio:} \texttt{"3331234567; 069876543"}
    \item \textbf{Vantaggi:} Non modifica la struttura del database e non richiede tabelle aggiuntive.
    \item \textbf{Svantaggi:} Impedisce al database di effettuare controlli sui singoli valori (es. non si può usare un vincolo \texttt{UNIQUE} per impedire che due utenti inseriscano lo stesso numero). Rendere aggiornabile o ricercabile un singolo numero richiede complesse operazioni di manipolazione delle stringhe.
\end{itemize}

\subsubsection*{2. Creazione di una Classe Esterna (Associazione 1:N)}
Si estrapola l'attributo trasformandolo in un'entità autonoma debole, legata all'entità principale da un'associazione uno a molti. È la soluzione relazionale per eccellenza.
\begin{itemize}
    \item \textbf{Schema:} 
    \begin{itemize}
        \item $\texttt{PERSONA}(\underline{\text{cf}}, \text{nome})$
        \item $\texttt{TELEFONO}(\underline{\text{numero}}, \underline{\underline{\text{cfPersona}}})$
    \end{itemize}
    \item \textbf{Vantaggi:} Scalabilità infinita (una persona può avere da zero a mille telefoni senza sprecare byte). Le query sono pulite e veloci, e si possono imporre vincoli di integrità sui singoli numeri.
    \item \textbf{Svantaggi:} Richiede un'operazione di \texttt{JOIN} per estrarre l'elenco completo dei telefoni associati a una persona.
\end{itemize}

\begin{tip}[Principi Generali di Scelta]
La scelta della tecnica di risoluzione è quasi sempre a favore dell'entità esterna, tranne in casi specifici di puro logging:
\begin{enumerate}
    \item \textbf{Dati operativi e ricercabili (Regola Aurea $\rightarrow$ Entità Esterna):} Se il dato ha valore di business, se il sistema dovrà mai cercare "a chi appartiene questo numero", se l'utente deve poter modificare o cancellare un singolo valore dalla sua lista, la creazione di una tabella esterna 1:N è l'unica via ingegneristicamente solida.
    \item \textbf{Dati descrittivi morti (Eccezione $\rightarrow$ Accorpamento Testuale):} Si opta per la stringa formattata solo per dati passivi, che vengono scritti una volta e letti in blocco unicamente per stamparli a schermo. Ad esempio: un attributo multivalore \texttt{note\_aggiuntive} o \texttt{tag\_ricerca} in sistemi molto semplici. 
\end{enumerate}
\end{tip}

\section{Codifica delle Gerarchie (Ereditarietà)}

Il modello relazionale è intrinsecamente "piatto" e non supporta nativamente il concetto orientato agli oggetti di ereditarietà (superclasse e sottoclassi). Pertanto, quando nel Class Diagram si incontra una gerarchia di generalizzazione/specializzazione, è necessario "smontarla" trasformandola in tabelle ordinarie. 

L'ingegneria del software offre tre strategie principali per affrontare questa traduzione, ciascuna con precisi impatti su prestazioni, memoria e integrità dei dati.

\paragraph{1. Partizionamento Orizzontale (Accorpamento verso il basso)}
Questa soluzione prevede l'eliminazione della superclasse e la creazione di una tabella per ciascuna sottoclasse. Ogni tabella figlia "eredita" fisicamente copiando al proprio interno tutti gli attributi della generalizzazione.

\begin{itemize}
    \item \textbf{Schema logico:} $\texttt{PERSONA(nome, cognome)}$ scompare. Si creano \texttt{STUDENTE}(\underline{\texttt{id}}, \texttt{nome}, \texttt{cognome}, \texttt{matricola}) e \texttt{DOCENTE}(\underline{\texttt{id}}, \texttt{nome}, \texttt{cognome}, \texttt{stipendio}).
    \item \textbf{Vantaggi:} Interrogare una specifica sottoclasse è estremamente veloce (nessun \texttt{JOIN}).
    \item \textbf{Svantaggi e Limiti:} Questa strada è matematicamente percorribile \textbf{solo se la generalizzazione è totale} (ovvero se non esistono istanze della superclasse che non appartengano a nessuna sottoclasse). Inoltre, se si vuole fare una query su tutte le persone (indipendentemente dal tipo), il database deve eseguire una pesante operazione di \texttt{UNION} tra le tabelle.
\end{itemize}

\paragraph{2. Partizionamento Singolo (Accorpamento verso l'alto)}
Questa soluzione, all'opposto, prevede l'eliminazione delle sottoclassi. Si crea un'unica, grande tabella (la superclasse) che ingloba tutti gli attributi di tutte le specializzazioni, introducendo un attributo \textbf{discriminante} (un enumeratore) per capire di che "tipo" sia la riga.

\begin{itemize}
    \item \textbf{Schema logico:} $\texttt{PERSONA}(\underline{\text{id}}, \text{tipoPersona}, \text{nome}, \text{cognome}, \text{matricola}, \text{stipendio})$
    \item \textbf{Vantaggi:} È sempre applicabile ed è la soluzione più performante in assoluto (niente \texttt{JOIN}, niente \texttt{UNION}). Molto amata nello sviluppo web moderno per la sua velocità in lettura.
    \item \textbf{Svantaggi (Problema dei vincoli persi):} Genera tabelle "sparse", piene di valori \texttt{NULL} (uno studente avrà lo stipendio \texttt{NULL}). Soprattutto, fa perdere il rigore delle associazioni originarie. 
\end{itemize}

\begin{example}[Recupero dei vincoli di consistenza]
Nel Class Diagram, supponiamo che l'associazione \textit{"Insegna"} leghi la classe \texttt{Corso} ESCLUSIVAMENTE alla sottoclasse \texttt{Docente}. 
Utilizzando l'accorpamento verso l'alto, la tabella \texttt{CORSO} dovrà avere una chiave esterna che punta all'unica tabella rimasta: \texttt{PERSONA}.
Questo permette a livello strutturale l'assurdità di far insegnare un corso a uno Studente. Per recuperare l'informazione persa, si \textbf{deve} imporre un vincolo applicativo o un \textit{CHECK constraint} SQL che assicuri che la chiave esterna punti solo a righe in cui l'attributo discriminante \texttt{tipoPersona = 'Docente'}.
\end{example}

\paragraph{3. Partizionamento Verticale (Traduzione in Associazioni 1:1)}
Questa strategia preserva la struttura del diagramma originale. Si crea una tabella per la superclasse (con i soli attributi comuni) e una tabella per ogni sottoclasse (con i soli attributi specifici). Le tabelle figlie vengono legate alla tabella padre tramite un'associazione 1:1 rigorosa, utilizzando la tecnica della \textbf{Chiave Primaria Condivisa}.

\begin{itemize}
    \item \textbf{Schema logico:} 
    \begin{itemize}
        \item $\texttt{PERSONA}(\underline{\text{idPersona}}, \text{nome}, \text{cognome})$
        \item $\texttt{STUDENTE}(\underline{\underline{\text{idPersona}}}, \text{matricola})$
        \item $\texttt{DOCENTE}(\underline{\underline{\text{idPersona}}}, \text{stipendio})$
    \end{itemize}
    \item \textbf{Vantaggi:} Rispetta la terza forma normale, non spreca spazio con i \texttt{NULL} e mantiene perfettamente validi i vincoli delle associazioni. Nel caso di \textbf{ereditarietà disgiunta}, è necessario aggiungere un \textit{trigger} che impedisca allo stesso \texttt{idPersona} di comparire sia in \texttt{STUDENTE} che in \texttt{DOCENTE}.
    \item \textbf{Svantaggi:} È la soluzione più costosa a livello computazionale. L'estrazione dei dati completi di un utente richiederà sempre un \texttt{JOIN} tra la tabella specifica e la superclasse.
\end{itemize}

\begin{tip}[Principi Generali di Scelta dell'Ereditarietà]
Non esiste una soluzione perfetta a priori, ma la pratica ingegneristica suggerisce queste euristiche:
\begin{enumerate}
    \item \textbf{Usa l'accorpamento in alto (Singola Tabella):} Quando le sottoclassi differiscono per pochissimi attributi e non hanno relazioni specifiche pesanti. Il piccolo spreco di spazio dovuto ai \texttt{NULL} è ampiamente ripagato dalla velocità estrema delle query.
    \item \textbf{Usa l'accorpamento in basso (Tabelle Separate):} Quando la generalizzazione è totale, le query non cercano quasi mai la "superclasse" generica, e le sottoclassi sono entità molto ricche e complesse di per sé.
    \item \textbf{Usa le associazioni 1:1 (Tabella per Classe):} Quando l'integrità dei dati è la priorità assoluta, le sottoclassi hanno tantissime associazioni specifiche con altre tabelle del database e lo spazio su disco deve essere ottimizzato al massimo.
\end{enumerate}
\end{tip}

\section{Riepilogo: Il Ciclo di Vita della Progettazione dei Dati}

Prima di introdurre il linguaggio di interrogazione e definizione dei dati, è utile schematizzare l'intera metodologia di progettazione di basi di dati affrontata finora. Il processo ingegneristico segue un flusso di raffinamento progressivo (dal livello astratto al livello implementativo), suddiviso in fasi distinte e sequenziali. Il quarto punto rappresenta un'anticipazione di quanto vedremo nel prossimo capitolo. 

\begin{enumerate}
    \item \textbf{Progettazione Concettuale (Modellazione del Dominio)}
    \begin{itemize}
        \item \textit{Obiettivo:} Rappresentare la semantica della realtà di interesse in modo puro e indipendente dalla tecnologia.
        \item \textit{Artefatti prodotti:} \textbf{Class Diagram UML} iniziale.
        \item \textit{Documentazione a corredo:} \textbf{Dizionari dei Dati} (dizionario delle classi, dizionario delle associazioni e dizionario dei vincoli di business non esprimibili graficamente).
    \end{itemize}

    \item \textbf{Ristrutturazione dello Schema (Revisione Concettuale)}
    \begin{itemize}
        \item \textit{Obiettivo:} Adattare il modello concettuale puro ai limiti e ai requisiti del modello logico di destinazione (es. eliminazione attributi strutturati/multivalore, gestione delle gerarchie, ottimizzazione delle chiavi).
        \item \textit{Artefatti prodotti:} \textbf{Class Diagram Revisionato}.
        \item \textit{Documentazione a corredo:} Aggiornamento dei dizionari. In particolare, il \textbf{Dizionario dei Vincoli} subisce un forte arricchimento per recuperare semanticamente tutte le informazioni perse durante gli appiattimenti strutturali (es. vincoli per le ereditarietà accorpate).
    \end{itemize}

    \item \textbf{Progettazione Logica (Traduzione Relazionale)}
    \begin{itemize}
        \item \textit{Obiettivo:} Tradurre gli elementi del diagramma revisionato in strutture dati compatibili con l'algebra relazionale.
        \item \textit{Artefatti prodotti:} \textbf{Schemi Relazionali} definitivi, con esplicita indicazione degli identificatori (\textit{Chiavi Primarie - PK}) e dei legami di navigabilità (\textit{Chiavi Esterne - FK}).
    \end{itemize}
    
    \item \textbf{Implementazione e Definizione dei Metadati (Verso il livello Fisico)}
    \begin{itemize}
        \item \textit{Obiettivo:} Istruire materialmente il motore del database (DBMS) a creare le strutture logiche appena progettate.
        \item \textit{Strumento:} Utilizzo del linguaggio \textbf{SQL} (in particolare il sotto-insieme DDL, \textit{Data Definition Language}), attraverso il quale gli schemi relazionali e i vincoli vengono codificati nei metadati del sistema.
    \end{itemize}
\end{enumerate}











